\title{Large Language Models Can Automatically Engineer Features for Few-Shot Tabular Learning}

\begin{abstract}
The exceptional capacity of Large Language Models (LLMs) to address complex and novel reasoning challenges positions them as promising tools for advancements in tabular learning—a field integral to various practical domains. We introduce \textsf{FeatLLM}, an innovative in-context learning approach where LLMs serve as automated feature creators, transforming raw input datasets to optimize tabular prediction tasks. These synthesized features are subsequently utilized by straightforward machine learning algorithms, such as linear regression, enabling efficient and high-performing few-shot learning. Importantly, at the inference phase, the \textsf{FeatLLM} paradigm relies exclusively on the identified features and a basic predictive model, removing dependence on ongoing LLM interactions. In contrast to other LLM-driven solutions, our method obviates the necessity of querying the LLM for each input during prediction and only demands access at the API level, circumventing typical limitations such as prompt size restrictions. Experimental results spanning various tabular datasets reveal that \textsf{FeatLLM} produces superior decision rules, achieving a notable average performance gain (10\%) over competing strategies like TabLLM and STUNT.
\end{abstract}

\section{Introduction}
Large Language Models (LLMs) have recently distinguished themselves through their capability to generate contextually apt outputs for new tasks, even in the absence of dedicated fine-tuning~\cite{creswell2022selection,imani2023mathprompter,taylor2022galactica}. Leveraging their extensive training on massive textual datasets, LLMs embed a broad spectrum of background knowledge that can replace domain experts in numerous settings~\cite{Genomes2021,singhal2023large,wilcox2003role}. This quality has propelled automation across disparate applications including dialogue interpretation~\cite{finch2023leveraging} and code correction~\cite{fan2023automated}. Critically, by incorporating concise task instructions and a small number of demonstrative cases into prompt designs, LLMs can infer the context and strategically focus on pivotal elements~\cite{brown2020language,zhao2023survey}. These characteristics underscore LLMs’ aptitude for data examination, enhancing their role as automated feature constructors.

The extension of LLMs’ underlying reasoning abilities and knowledge integration has recently permeated the field of tabular learning. Knowledge, such as patterns associating advanced age with greater susceptibility to disease, becomes valuable in clinical prediction tasks. Recent studies propose constructing natural language representations for tasks and features, subsequently converting the structured data into serialized text and submitting it as input to an LLM~\cite{dinh2022lift,wang2023anypredict}. Afterward, solutions commonly involve either in-context learning~\cite{nam2023semi} or efficient tuning strategies~\cite{dinh2022lift,hegselmann2023tabllm}. The generalization capacities derived from LLMs’ priors have been shown to outperform classic tabular learning techniques, particularly in scenarios where data is sparse~\cite{hegselmann2023tabllm}.

Existing LLM-driven tabular learning strategies exhibit several drawbacks. Firstly, when used for end-to-end predictions, these methods entail a minimum of one LLM inference per data instance, leading to substantial computational burdens. Many methods also demand fine-tuning of the LLM to attain desirable accuracy, restricting applicability to scenarios where training codebases or tuning APIs are unavailable—a pertinent issue considering that leading, state-of-the-art LLMs increasingly restrict access to inference-only APIs~\cite{anil2023palm,openai2023gpt4}. Additionally, most current frameworks perform poorly with extended prompts~\cite{liu2023lost}; as tabular datasets expand in feature count and consequently text serialization length, these techniques often become impractical or suffer diminished effectiveness. Collectively, these obstacles represent significant impediments for deploying LLM-based tabular learning systems in practical contexts.

We propose an alternative methodology that reframes the function of LLMs: rather than relying on them exclusively for end-to-end prediction, we task them with feature engineering, utilizing their pre-existing knowledge more efficiently. Concretely, our approach introduces a new in-context learning paradigm, \textsf{FeatLLM}, which seeks to elucidate the “criteria” guiding LLM-derived predictions. For example, in disease prediction tasks, the LLM can identify and articulate explicit rules specifying the conditions of features that signal the presence of a disease. These inferred criteria are subsequently used to construct novel features, offering replacements for the original features in downstream tasks. This paradigm positions the LLM as a generator of features, thereby simplifying subsequent modeling; after feature creation, an elaborate machine learning architecture is no longer required for inference, resulting in reduced inference delays. By harnessing the LLM's in-context learning via demonstrations inserted into prompts, features can be extracted without conventional training procedures. Moreover, adopting ensemble techniques such as feature bagging~\cite{breiman2001random} helps address issues related to large prompt sizes.

\textsf{FeatLLM} decomposes the task of engineering new features from prompts into two distinct stages: (1) analyzing the underlying problem to discern how features connect to targets; and (2) employing insights from this analysis together with a handful of illustrative examples to derive explicit, class-separating rules. By leading LLMs through this stepwise reasoning process, only salient features are retained for rule formulation. These resulting rules are translated into executable code and applied to the dataset, generating binary features that indicate whether each instance complies with the corresponding rule. A simple model is then trained on these binary features to provide class probability estimates. To bolster robustness, an ensemble strategy is adopted—features are repeatedly created with bagging over both features and samples. Carefully balancing the count of sampled features and instances during bagging also alleviates the issue of overly lengthy prompts. Since \textsf{FeatLLM} does not require training and keeps inference computationally light, it readily suits a variety of practical deployment scenarios.

Thirteen tabular datasets are used to benchmark \textsf{FeatLLM} under low-shot learning conditions, demonstrating consistently high and reliable effectiveness. Our results highlight the LLMs’ capability to blend prior information with dataset-derived patterns, facilitating the generation of valuable features. Across different experimental contexts, our approach surpasses prevailing few-shot learning baselines. An anonymized version of the implementation can be found at~\texttt{https://github.com/Sungwon-Han/FeatLLM}.

\section{Related Work}

\subsection{Few-Shot Learning with Tabular Data} Few-shot learning has revolutionized the ability to derive transferable features from datasets while minimizing labeling costs~\cite{chen2018closer}. Although early advancements were predominantly concentrated on image data, the applicability of few-shot paradigms now extends across multiple domains, such as text, audio, and notably, tabular formats~\cite{majumder2022few,nam2023stunt,schick2021exploiting}. The limitation of labeled samples in few-shot scenarios often increases susceptibility to learning misleading associations~\cite{bartlett2021deep}. To address these challenges, recent methods incorporate auxiliary information or exploit domain-specific prior knowledge in order to embed suitable inductive biases during training. Examples include the integration of substantial amounts of unlabeled tabular data, complemented by suitable augmentation strategies for semi-supervised learning~\cite{ucar2021subtab,yoon2020vime}, as well as leveraging unsupervised pre-training techniques that offer effective model initialization regardless of task~\cite{somepalli2021saint}. Unsupervised objectives can take the form of masking or intentional corruption of portions of input data, followed by reconstruction-based learning~\cite{arik2021tabnet,majmundar2022met}, or the application of contrastive learning through data augmentation~\cite{bahri2022scarf,chen2020simple}. Nevertheless, the intrinsic heterogeneity of tabular datasets complicates the identification of universal augmentation methodologies, and such approaches have shown limited effectiveness in few-shot contexts~\cite{nam2023stunt}.

More recently, research has shifted towards training models on a comprehensive assortment of tabular datasets that may be unrelated to the specific target task, before transferring these models onto the task of interest~\cite{zhang2023generative,levin2022transfer,zhu2023xtab}. TransTab~\cite{wang2022transtab}, for example, leverages transformer-based architectures to learn semantic associations among table columns via column names instead of relying solely on cell values. Insights drawn from numerous tables and column relationships enable zero-shot and few-shot predictions on previously unseen tabular tasks. In parallel, TabPFN~\cite{hollmann2022tabpfn} synthesizes datasets using mixtures of distinct distributions to pre-train a Bayesian neural network. Post pre-training, inference consists of directly inputting both training and test data from the target task, eliminating the need for additional training. This broad prior knowledge gathered from a spectrum of datasets has demonstrated superiority over conventional tree-based models and has exhibited clear benefits for few-shot performance.

\subsection{Language-Interfaced Tabular Learning} Utilizing prior knowledge through large language models (LLMs) constitutes a compelling strategy~\cite{anil2023palm,openai2023gpt4}. LLMs, which are trained on voluminous and varied text datasets, excel at generalizing to new tasks and have been successfully applied in several fields~\cite{brown2020language}. Recent works strive to tap into the vast and nuanced prior knowledge inherent in LLMs for enhancing tabular learning. Approaches commonly entail converting tabular information into textual representations and embedding them within LLM input prompts~\cite{dinh2022lift,hegselmann2023tabllm,wang2023anypredict}. By furnishing LLMs with prompts that encapsulate tabular data, detailed feature descriptions, and task-related information, they can capture complex task-feature relationships for inference purposes. Research directions encompass specialized model training with parameter-efficient adaptation strategies, such as LoRA~\cite{hu2021lora} and IA3~\cite{liu2022few}, or adopting in-context learning which simply augments prompts with a handful of demonstration examples without additional training~\cite{dinh2022lift,hegselmann2023tabllm,nam2023semi,slack2023tablet}.

Although previous strategies have demonstrated significant advancements in few-shot tabular learning, consistently relying on LLMs for inference imposes limitations, particularly in industrial settings. The present work seeks to move away from the standard black-box paradigm where every sample prediction is produced end-to-end through an LLM. Rather, it focuses the LLM's capacity on deducing the governing conditions or rules for each target class. Through \textsf{FeatLLM}, these rules are extracted and converted into executable program code, resulting in the construction of novel features. This approach enables predictions that bypass the LLM altogether, harnessing in-context learning to derive rules from prior knowledge and limited examples.

\section{Methods}
The core mechanism of \textsf{FeatLLM} centers on extracting ``rules"—class-specific conditions—from a limited number of input samples, as shown in Figure~\ref{fig:main_model}. These conditions are generated by an LLM, parsed, and subsequently transformed into binary features for each sample, which indicate whether a rule is met. The binary feature vectors are then subjected to a dot product with a set of non-negative trainable weights to quantify each class's probability. The training process tunes these weights, revealing the relative significance of each feature with respect to class prediction (see Section~\ref{Sec:training}). To increase robustness, the procedure is iterated using bagging, and the resulting models are aggregated via ensembling. Further elaboration on the problem definition and each methodological component is provided below.

\vspace*{-0.12in}
\paragraph{Problem formulation.} We analyze a tabular dataset composed of $N$ labeled instances, $\mathcal{D} = \{(\mathbf{x}^i, \mathbf{y}^i)\}_{i=1}^{N}$. Each instance $\mathbf{x}^i$ is a vector in $\mathbb{R}^d$, representing the $d$ input features of the dataset, and these features are accompanied by natural language names, $F = \{f_j\}_{j=1}^{d}$, with separate definitions supplied. Within the context of a classification problem, every $\mathbf{y}^i$ corresponds to a class drawn from a specified set. In the $k$-shot learning scenario, the model is trained using only $k$ ($<N$) labeled examples, which are randomly selected from the dataset.

\subsection{Prompt Design for Extracting Rules}
\label{Sec:prompt_design}
To enhance the accuracy of the LLM's reasoning in rule extraction, we structure the prompt to emulate the methodology an expert human would employ when tackling tabular data tasks. As outlined in Fig.~\ref{fig:prompt_for_rule} and Appendix~\ref{sec:example_rule}, the prompt is organized into three distinct sections:

\vspace*{-0.12in}
\paragraph{Basic information description.} This section delivers the necessary context to address the task (highlighted as orange in Fig.~\ref{fig:prompt_for_rule}). It covers the objective of the problem, including target label information, detailed descriptions of features, as well as illustration through sample instances. The task itself is articulated in a question format (such as ``Does this patient's myocardial infarction complications data suggest chronic heart failure? Yes or no?"). Additional examples are available in Appendix~\ref{sec:task_desc}. Feature descriptions specify the type of data for each variable, with optional inclusion of feature definitions. If a feature is categorical, sample category values can also be presented. For illustrative purposes, selected training examples are converted into textual format together with their actual labels. The approach for serialization follows practices established by previous works~\cite{dinh2022lift,hegselmann2023tabllm}:

\begin{align}
&\text{Serialize}(\mathbf{x}^i,\mathbf{y}^i, F) = \nonumber \\ 
&``\ f_1\ \text{is}\ \mathbf{x}^i_1.\ ... \ f_d\  \text{is}\  \mathbf{x}^i_d.\ \ \text{Answer:}\  \mathbf{y}^i\ ”
\end{align}

\vspace*{-0.12in}
\paragraph{Reasoning instruction.} Our approach seeks to improve LLM reasoning capabilities by imparting targeted instructions (refer to the blue text in Fig.~\ref{fig:prompt_for_rule}). The instruction sequence opens with an introductory statement modeled after the chain-of-thought methodology~\cite{wei2022chain}. Subsequently, we delineate the LLM's reasoning workflow into two distinct phases. In the initial phase, the model is prompted to hypothesize causal relationships or inclinations between individual features and the task description, drawing solely upon its general knowledge and intuition rather than referencing illustrative examples. This strategy enables the LLM to harness and articulate its foundational understanding of the problem space. The following phase utilizes both the exemplar demonstrations provided and the insights gleaned from the prior phase, guiding the LLM in deriving a set of rules corresponding to each class. This dual-phase process minimizes the risk of the model learning irrelevant correlations from non-informative columns and instead steers attention toward features of greater consequence. To maintain consistency and prevent both the proliferation of irrelevant rules and the loss of critical ones—which could result in either underfitting or an unwieldy prompt—we designate a constant number of rules (i.e., 10)\footnote{Discussion concerning rule quantity can be found in Section~\ref{sec:analysis}.} per class. Moreover, the rule construction process incorporates references to feature types described in the fundamental information section, mitigating potential mismatches between rule conditions and feature formats.

\vspace*{-0.12in}
\paragraph{Response instruction.} To streamline the extraction and subsequent application of generated rules, the LLM is systematically instructed to format its outputs according to a specified set of guidelines (refer to yellow text in Fig.~\ref{fig:prompt_for_rule}). The prompt communicates both the expected organizational framework for answers by class (i.e., response format) and the prescribed template for rules themselves (i.e., rule format). This facilitates the LLM's ability to match formatting standards with the feature type in question. In the absence of explicit directives, the LLM is capable of merging multiple rules by employing logical connectives such as AND or OR. Appendix~\ref{sec:outcome-rule} showcases an illustrative example of a completed outcome.

\subsection{Inferring Class Likelihood via Rules}
\label{Sec:training}

\paragraph{Parsing rules for feature generation.} In this phase, the previously identified rules are employed to generate novel binary features. These features are specific to each class, flagging whether a sample complies with the respective class's rule set, using a value of '0' or '1'. They are useful for estimating a sample's class membership probability. However, because rules derived from the LLM are in natural language form, a dedicated parsing method must be devised to automate feature extraction. Although formatting restrictions are imposed during the rule creation step to simplify parsing, outputs from the LLM may contain inconsistent syntactic variations~\cite{kovavc2023large}. For example, the model could replace ``$>$” or ``$<$” symbols with corresponding phrases like ``greater than'' or ``less than'', which complicates the parsing procedure. 

To mitigate issues posed by unstructured or noisy text, rather than designing intricate programmatic logic, our approach is to take advantage of the LLM's own capabilities. The LLM excels in interpreting semantic content irrespective of superficial textual modifications, and is proficient in writing concise code following detailed instructions due to its exposure to code-oriented datasets. Thus, we construct prompts that specify the function name, its input and output formats, as well as the relevant extracted rules, which are then provided to the LLM. Illustrative examples and outcomes for this prompting technique are documented in Figures~\ref{fig:prompt_for_parsing} and ~\ref{sec:example_parsing}-\ref{sec:example_parse_outcome} in the Appendix. The resulting Python code is executed via the \texttt{exec()} method to facilitate feature extraction.

\vspace*{-0.12in}
\paragraph{Inferring class likelihood.} With the binary features for each class derived from rule-based parsing, an elementary approach to gauge a sample’s likelihood of belonging to a particular class involves tallying the number of satisfied rules for each class—effectively computing the sum across that class’s binary features. However, the relative weight or informativeness of individual rules and features often varies, and it is important to determine these weights from training data. We propose using standard machine learning (ML) techniques to learn the contribution of each binary feature per class. For instance, a simple linear model without an intercept term can be applied. Define the binary feature vector produced for sample $\mathbf{x}^i$ and class $k$ as $\mathbf{z}_k^i \in \{0, 1\}^R$, where $R$ is the number of rules for class $k$, and $c$ is the total number of classes. The probability for each class $\mathbf{p}^i$ is evaluated using:

\begin{align}
\text{logit}_k^i &= \max(\mathbf{w}_k, 0) \cdot \mathbf{z}_k^i, \nonumber \\
\mathbf{p}^i &= \text{Softmax}({[\text{logit}_{1}^i, ..., \text{logit}_{c}^i]}). \label{eq:infer_prob}
\end{align}

In this context, $\mathbf{w}_k$ denotes the adjustable parameters in the linear framework, reflecting the significance attributed to each feature. By restricting the weights to a positive domain, it is guaranteed that features generated by the LLM are prevented from diminishing the probability assigned to any class during training. Subsequently, the resulting logit vector undergoes transformation into class probabilities using the softmax function. The learning of $\mathbf{w}_k$ is achieved by minimizing the cross-entropy loss with a limited set of training examples.

\vspace*{-0.12in}
\paragraph{Ensembling with bagging.} The procedure is iteratively applied, resulting in the generation of multiple inference models whose predictions are aggregated in a final ensemble. Specifically, by collecting the weights obtained from $T$ independent runs, represented as $\{\mathbf{w}[t]\}_{t=1}^T$, the ensemble prediction is calculated as the mean class probability $\mathbf{p}^i$ across all runs, utilizing the respective weight vectors $\mathbf{w}[t]$ per Eq.~\ref{eq:infer_prob}.

To foster variability in the rule extraction process for the ensemble, a high temperature is configured for LLM sampling, thereby increasing randomness. Additionally, the ordering of few-shot examples is shuffled, based on empirical findings that this ordering significantly affects LLM query results~\cite{lu2022fantastically}\footnote{Further discussion regarding rule diversity is presented in Appendix~\ref{sec:diversity}}. Bagging is incorporated to leverage this prediction diversity, whereby for each trial, only a subset of features or instances is utilized; this is especially advantageous with larger training datasets or feature sets. This ensemble approach brings forth two primary benefits. Firstly, even when some runs yield erroneous rules from the LLM, correctly generated rules in other runs can counterbalance these errors, strengthening the framework’s reliability. This principle leverages LLM self-consistency~\cite{wang2022self}, since rules repeatedly produced across several trials are statistically more likely to be valid. Secondly, the ensemble mechanism mitigates prompt size restrictions inherent to LLMs. When the tabular dataset exceeds the maximal input capacity, bagging serves to downscale the prompt, preserving robust model performance. Although a single rule set might fail to encapsulate the dependencies of every feature, the ensemble comprising a multitude of weak learners over multiple trials can collectively compensate for the deficiencies in individual trials' feature or instance coverage.

\section{Experiments}
We assess the effectiveness of \textsf{FeatLLM} across a range of tabular datasets, conducting in-depth analyses to illuminate the operational mechanisms of our framework. Additional experimental findings, such as the impact of using different LLMs and qualitative assessments, can be found in the Appendix due to space limitations.

\subsection{Performance Evaluation}
\paragraph{Datasets.}
The experimental suite covers 13 datasets used for both binary and multi-class classification tasks: (1) Adult~\cite{asuncion2007uci}, which involves predicting whether an individual’s annual income exceeds \$50,000; (2) Bank~\cite{moro2014data}, designed to forecast customer subscription to term deposits; (3) Blood~\cite{yeh2009knowledge}, which predicts if a donor will make future donations; (4) Car~\cite{kadra2021well}, focusing on evaluating vehicle quality; (5) Communities~\cite{misc_communities_and_crime_183}, aimed at predicting crime rates in specific locations; (6) Credit-g~\cite{kadra2021well}, which classifies individuals as presenting good or bad credit risks; (7) Diabetes\footnote{\texttt{kaggle.com/datasets/uciml/pima-indians-diabetes-database}}, which determines the likelihood of an individual having diabetes; (8) Heart\footnote{\texttt{kaggle.com/datasets/fedesoriano/heart-failure-prediction}}, for identifying cases of coronary artery disease; and (9) Myocardial~\cite{misc_myocardial_infarction_complications_579}, which predicts chronic heart failure status. 

Furthermore, we incorporate newly released datasets and synthetic ones not part of LLM pre-training and rarely reported in prior work: (10) Cultivars, targeting prediction of soybean grain yield\footnote{\texttt{archive.ics.uci.edu/dataset/913}}; (11) NHANES, for classifying individuals’ age group based on their record\footnote{\texttt{archive.ics.uci.edu/dataset/887}}; (12) Sequence-type, where the task is to distinguish whether a sequence belongs to arithmetic, geometric, Fibonacci, or Collatz types; (13) Solution-mix, aimed at assessing if the concentration in a blend from four solutions surpasses 0.5. The first two datasets were only added to the UCI Machine Learning Repository after September 2023, ensuring they were absent during the pre-training of the LLM API (gpt-3.5-turbo-0613) underpinning our model. The latter two are artificial datasets, guaranteeing no prior exposure to the LLM API before evaluation.

Detailed attributes regarding scale and complexity for each dataset are presented in Table~\ref{Tab:basic-desc}. The datasets feature explicit attribute names and explanations. Data such as `Communities' and `Myocardial', each containing over 100 features, are especially demanding for LLMs because of their constrained context window sizes.

\vspace*{-0.12in}
\paragraph{Baselines.}
Nine baseline methods are used for performance comparison. The first group consists of traditional approaches to tabular learning: (1) Logistic regression (LogReg), (2) XGBoost~\cite{chen2016xgboost}, and (3) RandomForest~\cite{ho1995random}. The following two methods rely on access to unlabeled data: (4) SCARF~\cite{bahri2022scarf}, and (5) STUNT~\cite{nam2023stunt}. It should be noted that acquiring extensive unlabeled datasets may not be feasible in practical scenarios. Another contemporary baseline is (6) TabPFN~\cite{hollmann2022tabpfn}, which uses artificially generated datasets with broad distributions for pretraining. The remaining three baselines engage LLMs: (7) In-context learning (In-context)~\cite{wei2022emergent}, (8) TABLET~\cite{slack2023tablet}, and (9) TabLLM~\cite{hegselmann2023tabllm}. In-context learning operates by incorporating several training examples directly within the input prompt, bypassing any parameter adjustment. TABLET expands on this approach, supplementing the input prompt with insights from an external classifier in the form of rule-sets and prototypes for improved inference accuracy. TabLLM applies a parameter-efficient tuning method such as IA3~\cite{liu2022few} to train the LLM with limited samples.

\vspace*{-0.12in}
\paragraph{Implementation details.}
\textsf{FeatLLM} is built upon GPT-3.5 as its foundational LLM, but its framework is flexible with respect to LLM selection (refer to Appendix~\ref{sec:backbones} for comparative results using alternative backbones). When performing LLM inference, a temperature of 0.5 and the API's standard top-p value of 1 are applied. We configure the ensemble size to 20 and set the rule extraction count at 10. The influence of hyper-parameters is further analyzed in Figure~\ref{fig:hyperparameter2} and detailed in Appendix~\ref{sec:hyper-parameter}.
For the linear model, we adopt the Adam optimizer at a learning rate of 0.01 and train for 200 epochs, employing $k$-fold cross-validation to select the most suitable epoch. For baseline reproduction, GPT-3.5 is used in in-context learning setups, while TabLLM leverages T0~\cite{sanh2022multitask} in accordance with its original protocol. Importantly, TabLLM enforces the use of LLMs with available public checkpoints, precluding the use of GPT-3.5. Parameter selection for classic machine learning algorithms (LogReg, XGBoost, RandomForest) is conducted via Grid-search integrated with $k$-fold cross-validation; the parameter $k$ is assigned as either 2 or 4 to ensure the training set covers every class at least once. For expanded implementation details, see Appendix~\ref{sec:baselines}.

\vspace*{-0.12in}
\paragraph{Results.}
Tables~\ref{tab:main1} and \ref{tab:main2} display the results of performance evaluations carried out over several datasets. Each experiment was performed three times, varying the random splits of training and test sets; we report the means and standard deviations of AUC scores. In comparison to various baselines, \textsf{FeatLLM} emerges as either the leading or second-best model across most scenarios. Remarkably, our method achieves similar effectiveness using only 4 training examples (shots) per class as classical tabular baselines that require as many as 32 shots (refer to Fig.~\ref{fig:compares}). Although STUNT and TabPFN show marked improvements for particular tasks, their overall advantage over traditional machine learning techniques remains modest. Additionally, LLM-based methods such as in-context learning, TABLET, and TabLLM struggled to process tabular data with high feature dimensionality. In contrast, our approach, incorporating both bagging and ensemble strategies, remains robust even for datasets containing a large number of features and continues to deliver strong performance. It is also worth mentioning that our ensemble and bagging methods can theoretically be applied to other LLM frameworks; nonetheless, doing so would entail twice as many inference passes per data point, significantly raising computational demands and limiting their practical use.

\subsection{Analysis \& Discussion}
\label{sec:analysis}
\paragraph{Ablation study.} To determine the effect of key components within the model, we conducted ablation studies involving: (1) \textsf{-Tuning}: removing the parameter tuning for the linear layer and, instead, simply aggregating binary features for logit computation; (2) \textsf{-Ensemble}: bypassing the ensemble method; (3) \textsf{-Description}: excluding the utilization of feature descriptions; and (4) \textsf{-Reasoning}: skipping Step-1 in the reasoning instruction during rule construction.

The impact of these component removals, as measured by changes in average AUC across all datasets, is presented in Table~\ref{tab:ablation}. Across all scenarios, omitting any component led to diminished model performance. Of particular note, the advantage of parameter tuning is increasingly pronounced as the number of shots rises, highlighting the growing value of rule importance estimation in data-rich settings. By contrast, including feature descriptions and reasoning steps delivers the largest gains when training data is limited, implying that drawing on LLMs' prior knowledge is particularly beneficial in few-shot regimes. The ensemble mechanism introduced by our framework steadily enhances performance under all experimental conditions.

\vspace*{-0.12in}
\paragraph{Training \& inference time comparison.} To evaluate computational efficiency, we compare both training and inference times across various methods using the Adult dataset, which includes 16 training shots. All experiments utilize a single A100 GPU by default, except for TabLLM, which requires four A100 GPUs for model parallelism. Table~\ref{tab:cost} presents a breakdown of the runtimes. Notably, \textsf{FeatLLM} demonstrates inference times that are on par with standard baseline approaches. Conversely, other LLM-based techniques—such as In-context learning, TABLET, and TabLLM—incur considerably greater inference times. For training, \textsf{FeatLLM} shows similar runtime costs relative to other LLM-driven models. It makes just 30 API requests, which are budget-friendly and amenable to parallelization.

\vspace*{-0.12in}
\paragraph{Quality of features generated by \textsf{FeatLLM}.}
To assess the informativeness of rule-based features produced by \textsf{FeatLLM}, we perform a straightforward evaluation on downstream tasks. In particular, we compute the average AUROC between these features and the corresponding target labels, and determine the proportion of features with AUROC scores above 0.5 within each dataset. 
Table~\ref{tab:quality} displays a summary of these findings. The results demonstrate that most of the generated rules are relevant to the prediction target and supply valuable information for solving the task (see the column ``Before tuning"). Additionally, when the linear model is tuned with the training data, features with minimal relevance (i.e., those whose weights fall below a set threshold) are systematically pruned (see the column ``After tuning"). The threshold is chosen as 0.1, based on the distribution of learned weights.

\vspace*{-0.12in}
\paragraph{Impact of introducing spurious correlations.} To evaluate the capability of different methods in discarding spurious and irrelevant correlations in settings with limited labeled data, we performed a dedicated analysis. Specifically, we employed the Heart dataset and augmented it by randomly adding features from the Adult dataset that bear no relevance to the Heart prediction task. We subsequently assessed how model performance shifts as we incrementally increase the count of such unrelated features. To ensure equitable comparison across approaches, we used a consistent 8-shot labeled sample configuration and omitted any unlabeled samples from the protocol, which excludes techniques such as SCARF and STUNT from the evaluation.

Figure~\ref{fig:spurious} presents the model performance alterations resulting from the presence of spurious correlations. Notably, \textsf{FeatLLM} was comparatively resilient to the addition of spurious features, displaying minimal decline in predictive performance versus other baselines. This robustness can be attributed to the framework’s reliance on prior task-feature knowledge during rule extraction, which constrains the generation of rules to relevant columns and suppresses those associated with irrelevant features. Consequently, rules involving noisy columns constituted only about 1-2\% of all rules produced. Nonetheless, we also note that \textsf{FeatLLM} is susceptible to picking up spurious associations and misconstruing causal links if the added columns originate from datasets closely aligned with the Heart task (see Appendix~\ref{sec:spurious-appendix}).

\vspace*{-0.12in}
\paragraph{Hyper-parameter analysis}
\textsf{FeatLLM} is governed by two primary hyper-parameters: the size of the ensemble and the number of rules generated per class with each LLM invocation. We ran experiments to systematically evaluate how these choices affect the overall performance. The findings indicate a trend of improved performance with a larger ensemble size; however, this improvement plateaus after approximately 20 ensembles, as illustrated in Fig.~\ref{fig:appendix-hyperparameter1} (Appendix). Regarding the rules per class, no significant statistical variations in performance were found, but selecting 10 rules appeared optimal in balancing prompt length against model efficacy (see Fig.~\ref{fig:hyperparameter2}). It is plausible that increasing the number of rules further inflates the input space—potentially augmenting the available information but also incurring a greater risk of overfitting, especially when sample sizes are limited.

\section{Conclusion}
We introduce \textsf{FeatLLM}, a method that elevates the performance ceiling for LLM-driven, few-shot tabular data learning. Distinct from conventional approaches that simply repurpose LLMs to classify or predict on a per-sample basis, \textsf{FeatLLM} leverages these models as feature engineers to derive informative prediction criteria. This is accomplished via an iterative rule generation step and an ensemble learning mechanism employing bagging, thus ensuring that only minimal inference costs are involved, no training process is required (API access suffices), and limitations arising from tabular feature dimensionality are effectively addressed. The empirical results establish \textsf{FeatLLM} as a leading strategy for leveraging LLMs on real-world tabular datasets where data efficiency is crucial.

Currently, \textsf{FeatLLM} targets scenarios characterized by scarce training data. Future research will extend its scope to datasets comprising larger sample sizes by developing new feature construction approaches, considering alternative feature forms in addition to rules, and advancing interpretability methods. Such improvements will facilitate broader applicability across diverse machine learning problems.

\section{Broader Impact}
The \textsf{FeatLLM} framework capitalizes on the inherent knowledge contained within LLMs to tackle tabular learning tasks, reaching levels of performance unattainable by standard supervised learning, particularly in regimes with limited data. By integrating pretrained domain-specific knowledge, the method helps address overfitting common in small data settings. This capability is especially relevant to practical contexts (e.g., finance, healthcare) where labels are costly or logistically challenging to obtain and the pretrained LLMs encompass rich, relevant textual corpora. As future work towards broader deployment, it is critical to investigate how social biases and misinformation present in LLM priors may shape model behavior, since these biases can directly—and at times predominantly—affect prediction outcomes, perhaps more so than the limited available labeled data.

\end{document}